% =====================================================================
%  Chương 6 — Kết luận
% =====================================================================
\chapter{Kết luận}
\label{ch:ketluan}

\section{Trả lời sáu hạng mục}

\begin{enumerate}
  \item \textbf{Cơ sở lý thuyết.} Mỗi khái niệm của mạng hồi quy (tensor chuỗi, Simple RNN, BPTT, clipping, tế bào LSTM và
        GRU, BiLSTM, recurrent dropout, LayerNorm) được viết thành một hàm NumPy và khớp với PyTorch ở cả
        \SL{theory/n_checks} phép đối chiếu. Phép đo gradient bằng autograd cho một kết quả đáng nhớ: LSTM chỉ giữ được
        gradient khi cổng quên mở. Với bias cổng quên bằng 0, gradient ở trễ 20 gần như bằng của Simple RNN.
  \item \textbf{Hai tập dữ liệu thực tế.} AMZN có \SL{amzn/n_raw} phiên, chia theo thời gian, dự báo return thay vì mức
        giá vì giá không dừng (ADF $p = \SL{amzn/adf/close/p}$). KKBox có \SL{kkbox/n_sample} người dùng thật lấy từ
        \SL{kkbox/log_rows} dòng log, đọc theo khối không cần nạp trọn vào RAM. Trên cả hai tập, mọi thống kê chuẩn hoá
        chỉ tính trên train.
  \item \textbf{PyTorch.} Bốn mô hình dùng chung một lớp và một vòng lặp huấn luyện, huấn luyện trên GPU sau khi đo được
        GPU nhanh hơn CPU. Trên KKBox, LSTM tốt nhất với F1 \SL{torch_kkbox/models/lstm/f1}\,\% và ROC-AUC
        \SL{torch_kkbox/models/lstm/roc_auc}; ba mô hình có cổng đều hơn Simple RNN.
  \item \textbf{Keras và đối chuẩn.} Cùng kiến trúc cho chất lượng gần như bằng nhau ở hai framework. PyTorch nhanh hơn
        với từng mẫu một (trung vị \SL{bench/lat_median/torch} so với \SL{bench/lat_median/keras} ms trên một luồng CPU),
        còn với lô lớn thì khoảng cách hẹp lại và Keras thắng ở Simple RNN. Tệp \texttt{.keras} nặng hơn vì chứa cả trạng thái
        optimizer.
  \item \textbf{Web App.} Flask phục vụ cả 16 mô hình qua REST API, giao diện hai tab; toàn bộ tham số tiền xử lý đọc từ
        \texttt{metadata.json} nên ứng dụng chạy được mà không cần dữ liệu gốc. Kịch bản kiểm chứng qua
        \SL{app/n_calls} lượt gọi.
  \item \textbf{Kết luận về mô hình.} Mạng hồi quy học được tín hiệu thật khi tín hiệu nằm trong chuỗi (xu hướng nghe nhạc
        giảm dần của người sắp rời bỏ). Nó không tạo ra được tín hiệu khi chuỗi không chứa (return cổ phiếu ngày mai):
        không mô hình nào trong 16 mô hình thắng mốc ``giá ngày mai bằng giá hôm nay''.
\end{enumerate}

\section{Hạn chế}

\begin{itemize}
  \item Mỗi mô hình chỉ chạy một seed. Chênh lệch giữa các mô hình trên KKBox (dưới hai điểm F1) và giữa hai framework
        đều nhỏ hơn mức dao động thường gặp giữa các lần khởi tạo, nên bài không xếp hạng chúng.
  \item KKBox chỉ dùng mẫu 100.000 người và 31 ngày log. Bỏ qua lịch sử giao dịch và thông tin thành viên nên ROC-AUC dừng
        ở khoảng 0,68, thấp hơn các lời giải đầy đủ của cuộc thi.
  \item AMZN chỉ dùng đặc trưng kỹ thuật tính từ chính chuỗi giá. Các sự kiện gây biến động lớn nhất (báo cáo quý, tin vĩ
        mô) không có trong dữ liệu.
  \item Keras chạy trên CPU vì TensorFlow không hỗ trợ GPU trên Windows, nên thời gian huấn luyện của hai framework không
        so trực tiếp được; bài chỉ so tốc độ suy luận trên cùng CPU.
\end{itemize}
